\section{Summary of Contributions}
\label{sec:conclusion_contributions}

\paragraph{Evaluation frameworks for open-world novelty.}
In \cref{ch:benchmarks}, we introduced \NovelGridworlds{} and \NovelGym{}, benchmark environments that enable the systematic study of agent behavior under novelty. These environments support controlled injection of diverse classes of mismatch---including changes to objects, environment dynamics, and task structure---organized according to the novelty taxonomy developed in \cref{sec:novelty_taxonomy}. They accommodate agents across multiple paradigms, including symbolic planning, reinforcement learning, and hybrid architectures, through flexible observation interfaces. They also provide metrics that capture the dynamics of adaptation: how sharply performance degrades following a change and how efficiently the agent recovers. These benchmarks established the experimental foundation for evaluating the subsequent contributions.

\paragraph{Action abstractions for efficient adaptation under novelty.}
In \cref{ch:rapid_learn}, we developed \RAPidLearn{}, a hybrid planning--learning framework that enables agents to recover from execution failures caused by novelty. The key insight is that structured action abstractions---the operator--executor decomposition formalized in \cref{sec:action_abstractions}---allow failures to be localized to specific components of a plan. When an executor fails because the environment dynamics have changed, the framework converts a global task failure into a targeted learning problem: the agent acquires only the new executor needed to restore successful execution while preserving previously valid behaviors. Experimental results demonstrated improvements of one to two orders of magnitude in sample efficiency over the evaluated reinforcement-learning and transfer-learning baselines. The neurosymbolic cognitive architecture further connected this mechanism to a cascading recovery pipeline, progressing from lightweight symbolic strategies to learning-based adaptation and evaluated across 288 Polycraft novelties. This contribution instantiates localized behavioral accommodation within the symbolic task interface studied in the dissertation.

\paragraph{Interaction abstractions for generalizable skill learning.}
In \cref{ch:flex,ch:pho}, we introduced object-centric, force-based representations for robotic manipulation that capture interaction at the level of object dynamics rather than robot-specific control. \FLEX{} demonstrated transfer to 13 previously unseen articulated objects and across three simulated robotic platforms without policy retraining, together with real-world validation on a physical UR5 robot. It also achieved more than an order-of-magnitude improvement in training efficiency over the evaluated robot-centric baselines. The pushing-heavy-objects framework extended the force-space approach to rigid-body transport under unknown mass and friction, demonstrating zero-shot deployment of a simulation-trained policy on a legged mobile manipulator. By defining policies in terms of forces applied to the object, both approaches capture task-relevant physical structure while abstracting away robot-specific control. They thereby expand the range of object, physical, and embodiment variation that can be assimilated within an existing policy representation.
